\documentclass{article}
\usepackage[utf8]{inputenc}
\usepackage{amsmath}
\usepackage[margin=1in]{geometry}
\usepackage{graphicx}
\usepackage{hyperref}

\title{ELCT 562 PA \#2}
\author{J. Vaught}
\date{\today}

\begin{document}

\maketitle

\section{Problem Statement and Data}
The objective is to analyze three different power delay profiles used in 4G LTE: EPA, EVA, and ETU. For each profile, we aim to determine:
\begin{itemize}
    \item Average Delay Spread
    \item RMS Delay Spread
    \item Coherence Bandwidth
\end{itemize}

The given data for each profile is as follows:

\textbf{EPA Delay Profile:}
\begin{verbatim}
Excess tap delay (ns)    Relative power (dB)
0                       0.0
30                      -1.0
70                      -2.0
90                      -3.0
110                     -8.0
190                     -17.2
410                     -20.8
\end{verbatim}

\textbf{EVA Delay Profile:}
\begin{verbatim}
Excess tap delay (ns)    Relative power (dB)
0                       0.0
30                      -1.5
150                     -1.4
310                     -3.6
370                     -0.6
710                     -9.1
1090                    -7.0
1730                    -12.0
2510                    -16.9
\end{verbatim}

\textbf{ETU Delay Profile:}
\begin{verbatim}
Excess tap delay (ns)    Relative power (dB)
0                       -1.0
50                      -1.0
120                     -1.0
200                     0.0
230                     0.0
500                     0.0
1600                    -3.0
2300                    -5.0
5000                    -7.0
\end{verbatim}

\subsection{General Solutions}

\subsubsection{Conversion to Linear Scale}
The relative power in dB (decibels) is converted to a linear scale to facilitate further calculations. The dB scale is logarithmic, and converting it to a linear scale makes it possible to perform arithmetic operations. The conversion formula is:
\begin{equation}
    \text{Power (linear)} = 10^{\frac{\text{Power (dB)}}{10}}
\end{equation}
This formula is derived from the definition of decibel, which is ten times the logarithm to base 10 of the ratio of the power to a reference power level.

\subsubsection{Average Delay Spread}
The average delay spread is calculated to assess the average extent of multipath delays. It represents the mean of the delays of different multipath components, weighted by their respective powers. The formula is:
\begin{equation}
    \text{Average Delay Spread} = \frac{\sum (\text{Delay} \times \text{Power})}{\sum \text{Power}}
\end{equation}
This is essentially the mean of the delay values, weighted by the linear power values.

\subsubsection{RMS Delay Spread}
The RMS (Root Mean Square) delay spread provides an estimate of the spread of multipath delays around the average delay spread. It is a measure of the width of the multipath delay profile. The formula is:
\begin{equation}
    \text{RMS Delay Spread} = \sqrt{\frac{\sum (\text{Power} \times (\text{Delay} - \text{Average Delay Spread})^2)}{\sum \text{Power}}}
\end{equation}
This formula calculates the square root of the average of the squared differences between each delay and the average delay, weighted by power.

\subsubsection{Coherence Bandwidth}
The coherence bandwidth is an estimate of the range of frequencies over which the channel can be considered “flat”. It is inversely related to the RMS delay spread. The formula used is an approximation:
\begin{equation}
    \text{Coherence Bandwidth} \approx \frac{1}{5 \times \text{RMS Delay Spread}}
\end{equation}
This approximation assumes that the channel is flat within the coherence bandwidth. The factor of 5 is a typical value used for calculating coherence bandwidth from the RMS delay spread.


\clearpage
\subsection{Specific Solutions for Each Profile}
\subsubsection{EPA Profile}
\textbf{Calculations and Results}
\begin{itemize}
    \item Average Delay Spread: 44.20 ns
    \item RMS Delay Spread: 43.13 ns
    \item Coherence Bandwidth: 4.64 MHz
\end{itemize}

\subsubsection{EVA Profile}
\textbf{Calculations and Results}
\begin{itemize}
    \item Average Delay Spread: 253.92 ns
    \item RMS Delay Spread: 356.65 ns
    \item Coherence Bandwidth: 0.56 MHz
\end{itemize}

\subsubsection{ETU Profile}
\textbf{Calculations and Results}
\begin{itemize}
    \item Average Delay Spread: 561.24 ns
    \item RMS Delay Spread: 990.94 ns
    \item Coherence Bandwidth: 0.20 MHz
\end{itemize}

\clearpage
\section{Calculation of Coherence Time}
Given the maximum Doppler spread frequency, \( f_{D} \), the coherence time, \( T_{c} \), can be estimated using the rule-of-thumb formula. The formula is given by:
\[ T_{c} \approx \frac{0.423}{f_{D}} \]
where \( T_{c} \) is the coherence time, and \( f_{D} \) is the maximum Doppler spread frequency.

Assuming that the maximum Doppler spread frequency is \( 50 \, \text{Hz} \), the coherence time can be calculated as follows:
\begin{align*}
T_{c} &= \frac{0.423}{50 \, \text{Hz}} \\
      &= 0.00846 \, \text{seconds} \\
      &= 8.46 \, \text{milliseconds}
\end{align*}

Therefore, the coherence time is approximately \( 8.46 \) milliseconds for a maximum Doppler spread frequency of \( 50 \, \text{Hz} \).

\clearpage
\section{Wireless Channel Analysis for Two Approaching Cars}

\subsection{Channel Fading Model}
Given the channel coherence bandwidth surpasses the signal bandwidth, the fading model is \textbf{flat-fading}. This follows because within the coherence bandwidth, phase and amplitude distortions affect all frequency components equally, leading to uniform fading across the signal spectrum.

\subsection{K-factor in Rician Model}
Considering the direct line-of-sight scenario of approaching cars, the \textbf{K-factor} would be large. This is characteristic of Rician fading where a predominant direct path, alongside fewer multipath components, skews the K-factor upwards, indicating a significant line-of-sight component.

\subsection{Carrier Frequency at the Receiver}
With both cars moving towards each other at \(10 \, \text{m/s}\) and a carrier frequency \( f_c = 3 \, \text{GHz} \), the received frequency, accounting for Doppler shift, is calculated as:
\[ f' = f_c \left(1 + \frac{v}{c}\right) \]
where \( v = 20 \, \text{m/s} \) (combined speed), and \( c \approx 3 \times 10^8 \, \text{m/s} \) (speed of light).

Performing the calculation:
\[ f' = 3 \times 10^9 \times \left(1 + \frac{20}{3 \times 10^8}\right) \approx 3,000,000,200 \, \text{Hz} \]
Thus, the received frequency is slightly higher than the transmitted 3 GHz due to the Doppler effect.

\end{document}
